\documentclass[10pt,compress,t,notes=noshow, xcolor=table]{beamer}

\input{../../style/preamble}
% Defines macros and environments
\input{../../latex-math/basic-ml.tex}
\input{../../latex-math/basic-math.tex}

\title{Interpretable Machine Learning}
% \author{LMU}
%\institute{\href{https://compstat-lmu.github.io/lecture_iml/}{compstat-lmu.github.io/lecture\_iml}}

\begin{document}

\titlemeta{
Interpretable Models 1
}{
Generalized Linear Models (GLMs)
}{
figure/logreg-2vars-surface
}{
\item Definition of GLMs
\item Logistic regression as example
\item Interpretation in logistic regression
}

%------------------------------------------------------------------
%------------------------------------------------------------------
\begin{framev}[align=top]{GLM \furtherreading{NELDERWEDDERBURN1972}}
\textbf{Problem}: Target variable given feat. not always normally distributed\\
$\leadsto$ LM not suitable
\begin{itemizeM}
\item Target is binary (e.g., disease classif.)\\
$\leadsto$ Bernoulli / Binomial distribution
\end{itemizeM}
\splitVCC[0.55]{
\begin{itemizeM}
\item Target is count variable\\
(e.g., number of sold products)\\
$\leadsto$ Poisson distribution
\item Time until an event occurs\\
(e.g., time until death)\\
$\leadsto$ Gamma distribution
\end{itemizeM}
}{
\image{figure/density_intervals.pdf}
\spacer[0]
}
\pause
\textbf{Solution}: GLMs extend LMs by allowing other distrib.-s from exp. family\\
\spacer[0.25]
\centerline{$g(\E(y \mid \xv)) = \xv^\top \thetav \; \Leftrightarrow \; \E(y \mid \xv) = g^{-1}(\xv^\top \thetav)$}
\spacer[0.25]
\begin{itemizeM}
\item Link function $g$ links linear predictor $\xv^\top \thetav$ to expectation of distrib. of $y \mid \xv$\\ %$\E$ of specified distribution of
$\leadsto$ LM is special case: Gaussian distrib. for $y \mid \xv$ with $g$ as identity func.
\item Link function $g$ and distribution need to be specified
\item High-order and interaction effects can be manually added as in LMs
\item Note: Interpretation of weights depend on link function and distribution
\end{itemizeM}
\end{framev}


\begin{framev}[align=top]{GLM - Logistic Regression}
\begin{itemizeM}
\item Logistic regression $\hat{=}$ GLM with Bernoulli distribution and logit link function:
$$
g(x) = \log\left(\frac{x}{1 - x}\right)
\Rightarrow \; g^{-1}(x) = \frac{1}{1 + \exp(-x)}
$$
\item Models probabilities for binary classification by
$$\pi(\xv) = \E(y \mid \xv) = P(y = 1) = g^{-1}(\xv^\top \thetav) = \frac{1}{1 + \exp(-\xv^\top \thetav)}$$
\end{itemizeM}
\centering
\image[0.3]{figure/logreg-2vars-surface.png}
\qquad
\image[0.4]{figure/reg_class_log_7}
\end{framev}


\begin{framev}[align=top]{GLM - Logistic Regression}
\begin{itemizeM}
\item Typically, we set the threshold to $0.5$ to predict classes, e.g.,
\begin{itemizeS}
\item Class 1 if $\pi(\xv) > 0.5$
\item Class 0 if $\pi(\xv) \leq 0.5$
\end{itemizeS}
\end{itemizeM}
\image{figure/reg_class_log_6.pdf}
% figs from i2ml course, Chapter 03.04: Logistic Regression, slide 7 
\end{framev}


%------------------------------------------------------------------
%------------------------------------------------------------------

\begin{framev}[align=top]{GLM - Log. Regression - Interpretation}
\begin{itemizeM}
\item \textbf{Recall:} Odds is ratio of two probabilities, odds ratio is ratio of two odds
\item Weights $\theta_j$ are interpreted linear as in LM (but w.r.t. log-odds)\\
$\leadsto$ difficult to comprehend
$$log\text{-}odds = \log\left(\frac{\pi(\xv)}{1 - \pi(\xv)}\right) = \log\left(\frac{P(y = 1)}{P(y = 0)}\right) = \theta_0 + \theta_1 x_1 + \ldots + \theta_p x_p$$
\textbf{Interpretation}: Changing $x_j$ by one unit changes log-odds of class 1 compared to class 0 by $\theta_j$
\pause
\item Odds for cls 1 vs. cls 0: %Odds can be caulculated by:
% deliberate choice to use in line with \displaystyle given full slide
$odds = \displaystyle\frac{\pi(\xv)}{1 - \pi(\xv)}= \exp(\theta_0 + \theta_1 x_1 + \ldots + \theta_p x_p)$
\item Instead of interpreting changes w.r.t. log-odds, it is more common to use the odds ratio
$$\frac{odds_{x_j+1}}{odds} = \frac{\exp(\theta_0 + \theta_1 x_1 + \ldots + \theta_j(x_j + 1) + \ldots + \theta_p x_p)}{\exp(\theta_0 + \theta_1 x_1 + \ldots + \theta_j x_j + \ldots + \theta_p x_p)} = \exp(\theta_j)$$
\textbf{Interpretation}: Changing $x_j$ by one unit changes the \textbf{odds ratio} for class 1 (compared to class 0) by the \textbf{factor} $\exp(\theta_j)$
%\pause
\end{itemizeM}
\end{framev}


\begin{framev}[align=top]{GLM - Logistic Regression - Example}
\begin{itemizeM}
\item Create a binary target variable for bike rental data:
\begin{itemizeS}
\item Class 1: ``high number of rentals'' $> 70\%$ quantile (i.e., \code{cnt} $> 5531$)
\item Class 0: ``low to medium number of rentals'' (i.e., \code{cnt} $\leq 5531$)
\end{itemizeS}
\item Fit a logistic regression model (GLM with Bernoulli distri. and logit link)
\end{itemizeM}
\splitVTT[0.55]{
\spacer[0]
\centering
%\begin{table}[ht]
% \begin{tabular}{rrrrr}
%   \hline
%  & Weights & SE & z value & p-value \\ 
%   \hline
% (Intercept) & -8.5 & 1.2 & -7.1 & 0.00 \\ 
%   seasonSPRING & 1.7 & 0.6 & 2.9 & 0.00 \\ 
%   seasonSUMMER & -0.9 & 0.8 & -1.1 & 0.26 \\ 
%   seasonFALL & -0.6 & 0.6 & -1.2 & 0.25 \\ 
%   temp & 0.3 & 0.0 & 7.4 & 0.00 \\ 
%   hum & -0.1 & 0.0 & -5.0 & 0.00 \\ 
%   windspeed & -0.1 & 0.0 & -3.0 & 0.00 \\ 
%   days\_since\_2011 & 0.0 & 0.0 & 11.6 & 0.00 \\ 
%   \hline
% \end{tabular}
\begin{scriptsize}
\begin{tabular}{rrrr}
\hline
& Weights & SE & p-value \\
\hline
(Intercept) & -8.52 & 1.21 & 0.00 \\ 
seasonSPRING & 1.74 & 0.60 & 0.00 \\ 
seasonSUMMER & -0.86 & 0.77 & 0.26 \\ 
seasonFALL & -0.64 & 0.55 & 0.25 \\ 
temp & 0.29 & 0.04 & 0.00 \\ 
hum & -0.06 & 0.01 & 0.00 \\ 
windspeed & -0.09 & 0.03 & 0.00 \\ 
days\_since\_2011 & 0.02 & 0.00 & 0.00 \\ 
\hline
\end{tabular}
\end{scriptsize}
%\end{table}
\pause
\spacer[0]
}{
\spacer[0]
%\pause
\imageC[0.9]{figure/logistic_marginal_temp.pdf}
\spacer[0]
}
\spacer[0.25]
\textbf{Interpretation}
\begin{itemizeM}
%\item If the temperature increases by 1 degree Celsius then the log odds of high number of bike rentals increase linearly by 0.3 c.p.
\item If \code{temp} increases by $1^\circ C$, odds ratio for class 1 increases by factor $\exp(0.29) = 1.34$ compared to class 0, c.p. ($\hat{=}$ ``high number of bike rentals'' now 1.34 times more likely)
%\item If \code{season} is SPRING instead of WINTER (reference), odds ratio for class 1 increases by factor $\exp (1.74) = 1.34$ compared to class 0
\end{itemizeM}
\end{framev}




% % --- Begin Embedded Slides on Marginal Effects in Logistic Regression ---

% \begin{frame}{Logistic Regression: Goal of Marginal Effect Calculation}
% \textbf{Objective:} Compute the instantaneous change in the predicted probability $p$ when the covariate $x_1$ (e.g., temperature) increases, holding other variables constant.

% \vspace{1ex}
% Model specification:
% \[
% p = \sigma(\eta), \quad \sigma(\eta) = \frac{1}{1 + e^{-\eta}}, \quad \eta = \beta_0 + \delta + \beta_1 x_1,
% \]
% where $\delta$ collects all other fixed covariate contributions.

% We are interested in:
% \[
% \frac{\partial p}{\partial x_1}.
% \]
% \end{frame}

% \begin{frame}{Marginal Effect: Chain Rule Decomposition}
% Apply the multivariable chain rule:
% \[
% \frac{\partial p}{\partial x_1}
%   = \frac{\partial p}{\partial \eta} \cdot \frac{\partial \eta}{\partial x_1}.
% \]

% \begin{itemize}
%   \item The first term captures the sensitivity of the sigmoid function.
%   \item The second term is the direct effect of $x_1$ on the linear predictor.
% \end{itemize}
% \end{frame}

% \begin{frame}{Sigmoid Derivative: Algebraic Derivation}
% Start with the definition:
% \[
% \sigma(\eta) = \frac{1}{1 + e^{-\eta}} = (1 + e^{-\eta})^{-1}.
% \]

% Differentiate:
% \[
% \frac{d\sigma}{d\eta} = - (1 + e^{-\eta})^{-2} \cdot (-e^{-\eta}) = \frac{e^{-\eta}}{(1 + e^{-\eta})^2}.
% \]

% Express in terms of $\sigma(\eta)$:
% \[
% \sigma(\eta) = \frac{1}{1 + e^{-\eta}}, \quad 1 - \sigma(\eta) = \frac{e^{-\eta}}{1 + e^{-\eta}},
% \]
% so:
% \[
% \frac{d\sigma}{d\eta} = \sigma(\eta)(1 - \sigma(\eta)) = p(1 - p).
% \]
% \end{frame}

% \begin{frame}{Linear Predictor Derivative with Respect to $x_1$}
% Assume:
% \[
% \eta = \beta_0 + \delta + \beta_1 x_1.
% \]

% Then:
% \[
% \frac{\partial \eta}{\partial x_1} = \beta_1.
% \]
% This captures the direct (linear) contribution of $x_1$ to $\eta$.
% \end{frame}

% \begin{frame}{Marginal Effect Formula for Covariate $x_1$}
% Combine results from the sigmoid and predictor derivatives:
% \[
% \frac{\partial p}{\partial x_1}
%   = \frac{\partial p}{\partial \eta} \cdot \frac{\partial \eta}{\partial x_1}
%   = p(1 - p)\,\beta_1.
% \]

% \begin{block}{Final Expression}
% \[
% \boxed{\displaystyle \frac{\partial p}{\partial x_1} = p(1 - p)\,\beta_1}
% \]
% \end{block}

% \textbf{Interpretation of the Marginal Effect}
% \begin{itemize}
%   \item $p(1 - p)$ is the slope of the sigmoid function:
%   \begin{itemize}
%     \item Maximal at $p = 0.5$ (most uncertain region).
%     \item Approaches $0$ as $p \to 0$ or $p \to 1$ (saturation).
%   \end{itemize}
%   \item $\beta_1$ scales the effect: larger values imply steeper increases in predicted probability with $x_1$.
%   \item The overall marginal effect is context-dependent and nonlinear in $x_1$ due to the form of $\sigma(\eta)$.
% \end{itemize}
% \end{frame}

% % --- End Embedded Slides on Marginal Effects ---






%------------------------------------------------------------------
%------------------------------------------------------------------

\endlecture
\end{document}
